% SHARD-ID: AQM-76-QSA-COTANGENT-FIRST-ASSOCIATION
% SHARD-TITLE: Cotangent-Based First Association
% SHARD-SUMMARY: Places the tangent-cotangent Weyl workflow into the QSA catalogue as a geometric first-association workflow.
% SHARD-SUMMARY: Derives the finite phase-label space from a point convention, cotangent fibre, tangent dual, and canonical symplectic pairing.
% SHARD-SUMMARY: Records the existing AQM-58--AQM-63 examples and leaves intrinsic-versus-relative and finite-ring database rows as gaps.
% SHARD-KEYWORDS: quantum-system association, cotangent, tangent, Weyl, first association, finite-residue points, kinematics

\section{Cotangent-Based First Association}
\label{sec:qsa-cotangent-first-association}

\subsection{Status and Local Anchors}

This is QSA Step 12 from
\path{docs/quantum_system_associations/PLAN.md}.  It is a new
synthesis/derivation shard: it adds no external source, no run artifact, and
no finite-ring database claim.  Its role is to place the already developed
tangent-cotangent Weyl workflow into the QSA catalogue as a geometric
first-association workflow.

The governing comparison language is AQM-65 and \texttt{CONVENTIONS.md}
convention \texttt{(ap)}.  The point-selection discipline is AQM-75.  The
technical convention is \texttt{CONVENTIONS.md} \texttt{(ao)}, and the local
technical anchors are AQM-58 through AQM-63.  The K\"ahler differential and
tangent-space source locators are inherited from AQM-58--AQM-61; this shard
does not introduce a new literature layer.

\subsection{Workflow Slot in the QSA Catalogue}

The cotangent-based first association is the following geometric workflow.
It starts from a geometric presentation, not from a bare finite set:
\[
  X=\Spec A
\]
over a named base field or base ring already fixed for the differential
calculation.  It then chooses a point convention in the sense of AQM-75.  For
the finite Weyl part of this shard, the selected points must be
finite-residue points; rational supports are allowed only when they are
explicitly named finite truncations, and closed finite-residue supports are
the current full finite-site convention for affine finite-type
\(\mathbb F_q\)-schemes.

For each selected finite-residue point \(x\), put
\[
  K_x=\kappa(x).
\]
Convention \texttt{(ao)} attaches the cotangent fibre
\[
  V_x^*=T^*_{X/k,x}=\Omega^1_{A/k}\otimes_A K_x
\]
and the tangent space
\[
  V_x=T_{X/k,x}=\operatorname{Hom}_{K_x}(V_x^*,K_x).
\]
The geometric phase-label space is
\[
  E_x^{\rm geom}=V_x\oplus V_x^*.
\]
The symplectic form is the canonical tangent-cotangent pairing
\[
  \omega_x((v,\alpha),(w,\beta))=\beta(v)-\alpha(w).
\]
This is the required symplectic datum before any Weyl representation is used.
It is bilinear and alternating by direct cancellation.  It is nondegenerate
when \(V_x\) is finite-dimensional over \(K_x\): a nonzero tangent vector is
detected by some cotangent vector, and a nonzero cotangent vector is detected
by some tangent vector.  This is the local derivation already recorded in
AQM-60.

If \(K_x\) is finite and \(m_x=\dim_{K_x}V_x<\infty\), convention
\texttt{(ao)} applies the finite Weyl-Heisenberg step to
\[
  E_x^{\rm geom}
\]
using the additive character
\[
  \psi_{K_x}(a)=
  \exp\!\left(\frac{2\pi i}{p}\operatorname{Tr}_{K_x/\mathbb F_p}(a)\right).
\]
The resulting kinematical output is a projective representation
\[
  \rho_x^{\rm geom}:E_x^{\rm geom}\to PU(\mathcal H_x^{\rm geom})
\]
with carrier dimension
\[
  \dim_{\mathbb C}\mathcal H_x^{\rm geom}=|K_x|^{m_x}.
\]
For a finite support \(S\), the phase labels and representations combine as
\[
  E_S^{\rm geom}=\bigoplus_{x\in S}E_x^{\rm geom},\qquad
  \rho_S^{\rm geom}=\bigotimes_{x\in S}\rho_x^{\rm geom}.
\]

Thus the first association output is not a direct assignment of a Hilbert
space to the set of points.  It is the sequence
\[
  \begin{gathered}
  \text{chosen points}\quad\leadsto\quad
  \text{cotangent/tangent fibres}\quad\leadsto\quad
  \text{symplectic phase labels}\\
  \leadsto\quad
  \text{Weyl-Heisenberg carriers when finite}.
  \end{gathered}
\]

\subsection{Relation to Earlier Association Families}

This workflow is different from the structureless finite-set first
association of AQM-67.  AQM-67 starts with a bare finite set and produces the
formal carrier \(\mathbb C\langle X\rangle\), with any Gram matrix marked as
extra data.  Here the input includes a scheme presentation, a base for
differentials, and a point convention.  The output is controlled by
\(\Omega^1_{A/k}\), not by cardinality alone.

It is also different from the residue-Weyl field-label workflow reviewed in
AQM-71.  The residue-Weyl layer uses the residue field shape
\[
  K_x^{(q)}\oplus K_x^{(p)}
\]
at a site.  The cotangent-first layer uses
\[
  T_{X/k,x}\oplus T^*_{X/k,x}.
\]
Convention \texttt{(ao)} records that these agree only after extra choices in
special cases, for example at a smooth curve point where
\(\dim_{K_x}T_{X/k,x}=1\) and a tangent basis identifies the geometric label
space with \(K_x^2\).  Product rings, higher-dimensional smooth points,
singular points, and nonreduced points show that the two workflows are not
the same catalogue entry.

\subsection{Anchored Example Inventory}

The existing shards supply the following checked or locally derived examples.
This table is an inventory of local anchors, not a finite-ring-wide cotangent
database.

\begin{center}
\small
\begin{tabular}{p{0.23\linewidth}p{0.27\linewidth}p{0.36\linewidth}}
\toprule
Object and point choice & Local cotangent-first output & Anchor and status \\
\midrule
\(k^n\) with its \(n\) spectrum points &
\(\Omega^1_{k^n/k}=0\), so every
\(E_x^{\rm geom}=0\) and every geometric carrier is one-dimensional. &
AQM-58 proves the vanishing by derivations; AQM-60 and AQM-61 record the
trivial tangent-Weyl reading.  This is a geometric zero case, not the
residue-qudit system of AQM-40. \\
\addlinespace
\(\Spec\mathbb F_3\) &
One point, zero cotangent and tangent spaces, and no nontrivial geometric
Weyl label. &
AQM-61 Example 1. \\
\addlinespace
\(\Spec\mathbb F_3[\epsilon]/(\epsilon^2)\) at its closed point &
One tangent direction and one cotangent direction; the finite carrier has
dimension \(3\). &
AQM-61 Example 3.  This is first-derivative geometry, not the Artin-ring
thickened Weyl layer. \\
\addlinespace
\(\Spec\mathbb F_3[t]/(t^3)\) at its closed point &
One tangent direction in characteristic \(3\); the carrier has dimension
\(3\). &
AQM-61 Example 4.  The first tangent dimension does not record the whole
length-three thickening. \\
\addlinespace
\(\Spec\mathbb F_3[u,v]/(u^2,v^2)\) at its closed point &
Two tangent directions; the carrier has dimension \(3^2=9\). &
AQM-61 Example 5. \\
\addlinespace
\(\mathbb A^1_{\mathbb F_3}\), rational support &
Each rational point has one tangent direction; the support carrier has
dimension \(3^3\). &
AQM-62 Example 6.  Closed degree-\(d\) points have residue field
\(\mathbb F_{3^d}\) and carrier dimension \(3^d\). \\
\addlinespace
\(\mathbb A^2_{\mathbb F_3}\), rational support &
Each rational point has two tangent directions; the support carrier has
dimension \(3^{18}\). &
AQM-62 Example 7. \\
\addlinespace
The crossing \(xy=0\) over \(\mathbb F_3\), rational support &
The origin has two tangent directions and the four other rational points have
one; the support carrier has dimension \(3^6\). &
AQM-62 Example 8. \\
\addlinespace
Finite-field circle \(x^2+y^2=1\) &
Odd-characteristic smooth points have one tangent direction; the
\(\mathbb F_3\) rational support carrier has dimension \(3^4\), and the
\(\mathbb F_5\) rational support carrier has dimension \(5^4\). &
AQM-59 computes the cotangent fibres and point counts; AQM-62 Examples 9--10
record the Weyl carrier dimensions.  AQM-59 also records the
characteristic-two nonreduced tangent thickening. \\
\bottomrule
\end{tabular}
\end{center}

\subsection{What AQM-63 Adds and What It Does Not Add}

AQM-63 records a compatible cotangent-label selection test:
\[
  L_{\rm dR}(S)=\operatorname{ev}^1_S(dA)
  \subset 0\oplus\bigoplus_{x\in S}T^*_{X/k,x}.
\]
Because the cotangent summand is isotropic, these labels define commuting
Weyl operators after phases are fixed.  This is a local compatibility result
inside the cotangent-first kinematics.

For QSA Step 12, this does not turn the workflow into dynamics.  AQM-63 also
records that tangent-type labels, Hamiltonians, action functionals, metrics,
Hodge identifications, and other dynamical or CSS-completion choices remain
extra data.  Therefore the present shard uses AQM-63 only to say that exact
one-forms can be read as optional cotangent/momentum constraints inside the
already built phase-label space.

\subsection{Gaps Carried to Later Steps}

The intrinsic-versus-relative cotangent problem is deliberately not settled
here.  Every displayed \(T^*_{X/k,x}\) uses the base appearing in convention
\texttt{(ao)} and in the anchor shard.  Step 13 must compare local/intrinsic
and relative cotangent choices for \(\Spec R\) with the base ring named in
each example.

No finite-ring-wide cotangent table is asserted.  In particular, this shard
does not claim populated database-backed cotangent rows, helper coverage, or
finite-ring example-matrix completeness.  Those rows remain gaps pending the
separate finite-ring cotangent work item \texttt{aqm-qsa-frcot01} and the
later QSA finite-ring matrix steps.

Finally, AQM-75 still controls point selection.  Generic points, infinite
residue fields, analytic completions, and all-scheme-point indexing do not
become finite Weyl systems by this shard.  If \(K_x\) is not finite, or if
the representation category is not fixed, the output stops at the geometric
phase-label space and the missing representation choice must be labelled as a
gap.
